\subsection{Integration of Prometheus}  
In the research on which database would be suitable for centralized data storage and analysis, Prometheus played a central role. Not only is it an open-source project \cite{prometheus_wiki}, but it is also well known in the industry as a reliable monitoring solution for complex systems. It has a large community, along with comprehensive documentation and strong support, which are essential for maintaining a stable IT environment.\\
Another reason for choosing Prometheus is its seamless integration with Grafana for visualization, as well as its powerful PromQL query language. Additionally, Prometheus can store metrics in a time series database, which is particularly important for the detection systems \cite{prom_advantages}.\\

Prometheus consists of several components that are required for its operation. The core of the system is the Prometheus server, which collects and stores metrics and provides access to them through the Prometheus Query Language (PromQL).\\
Client libraries are directly integrated into the application code and expose a client-server interface. The Prometheus server periodically fetches data from them, by default every 15 seconds, this interval is configurable \cite{Prometheus_how_it_works}.\\

As explained in the previous subsection, Prometheus collects data via metrics exporters \ref{lst:perf-prometheus-exporter}. After the data is collected and stored centrally in the Prometheus time series database, a Python script was developed that allows exporting the data within a defined time range and with specified time steps.\\
This export functionality is essential for enabling further analysis, including exploratory data analysis, using external tools such as Python or statistical software. The following paragraph presents the Python script developed to export Prometheus data into a structured \texttt{CSV} format, enabling its direct use in subsequent analytical workflows.\\

The Python script is designed to systematically retrieve and organize time-series metrics from a Prometheus monitoring server, facilitating subsequent analysis and interpretation. The workflow begins by defining a specific time range for the data of interest and partitioning this interval into smaller, manageable chunks, which helps prevent overwhelming the server with excessively large queries and ensures efficient data retrieval. For each feature specified in an external configuration file, the script constructs appropriate Prometheus queries, attempting first to obtain the standard metric and, if unavailable, a performance-prefixed variant. The retrieved data are then parsed and organized according to their respective timestamps, forming a structured representation of the system's behavior over time. Finally, the script outputs these metrics in a wide-format \texttt{CSV} file, where each row corresponds to a timestamp and each column corresponds to a specific metric. This approach not only allows for straightforward integration with downstream data analysis or visualization workflows but also ensures reproducibility, scalability, and clarity in the presentation of high-dimensional monitoring data.
